\section{Lensing and Cosmological Viability}

The galaxy-domain results of \S\S4--6 establish that the response sector works empirically for rotation curves. A credible $\Lambda$CDM-alternative candidate must additionally address lensing and cosmology. This section provides the first lensing-facing statement and the minimal cosmological viability sketch.

\subsection{Lensing: The Weyl Potential Constraint}

\begin{sloppypar}
Galaxy--galaxy lensing and cluster lensing constrain the Weyl potential \textbf{(see PDF; symbol/definition not recovered from DOCX extraction)}. In standard GR with a perfect-fluid source (including CDM), the gravitational slip
\textbf{(see PDF; slip definition not recovered)}
at leading order in the quasi-static weak-field limit. The response-sector framework must make a definite prediction for $\eta$.
\end{sloppypar}

\subsubsection{Working Hypothesis: Metric Equivalence}

As a first working hypothesis, we assume the response sector produces \textbf{(see PDF; ``zero slip'' condition not recovered)} (zero slip) in the galaxy quasi-static regime. This is the most $\Lambda$CDM-like option and the easiest to test: it predicts that the lensing mass inferred from the response-sector's rotation-curve fit matches the lensing mass measured directly. Specifically, the implied convergence $\kappa(R)$ can be computed from the fitted potential and compared against weak-lensing measurements.

\subsubsection{Gravitational Slip as a Discriminant}

If the response sector includes anisotropic stress \textbf{(see PDF; symbol/definition not recovered)}, it generically produces gravitational slip \textbf{(see PDF; symbol/definition not recovered)}. This would be a distinctive observational signature. We do not introduce a free slip function; rather, we note that the choice of microphysical completion (\S3.5) determines the slip prediction. If future lensing data require slip, this constrains the completion; if they exclude it, the completion must produce \textbf{(see PDF; no-slip condition)}.

\subsection{Cosmological Background}

For the response sector to serve as a CDM replacement on cosmological scales, its effective energy density must evolve approximately as \textbf{(see PDF; scaling not recovered)} (dust-like) over the redshift range relevant for structure formation and distance measurements. This is the background viability requirement: the response sector must not ruin the distance--redshift relation that $\Lambda$CDM fits well.

\subsection{Cosmological Perturbations}

At the perturbation level, the response sector must cluster appropriately to produce the observed matter power spectrum and CMB angular power spectrum. A dust-like clustering component with an effective sound speed close to zero \textbf{(see PDF; symbol/definition not recovered)} over the relevant scales would most closely mimic CDM perturbations. If the response sector introduces a significant effective sound speed or anisotropic stress at cosmological scales, it will modify the growth rate and ISW effect in potentially detectable ways.

\subsection{Domain Partition and Honest Scope}

We adopt explicit domain partitioning: this work establishes the galaxy-domain closure and demonstrates empirical viability at rotation-curve scales. The lensing and cosmological domains are enumerated as defined requirements, not claimed results. The language discipline is: ``This work establishes the galaxy-domain closure and enumerates the cosmological completion requirements and discriminants.''
